\chapterimage{health.pdf}
\chapterimagecaption{\textit{Lunch atop a Skyscraper (1932)} by Charles Clyde Ebbets?}

\chapter{Health and Safety}

    \section{Health Insurance} \index{Health Insurance}
    
\subsection{Canadian Students}

\subsubsection{PGSS Health Insurance}
    Canadian Graduate students are offered Health and Dental plans by PGSS. The billing for these is included in your student fees. The plans are separate, you can choose either one or the other, both of them, or neither. They include coverage on:
    \begin{itemize}
        \item Health (prescription drugs, vaccinations, discounts on physio, massage, chiro, etc.)
        \item Vision (glasses, some eye exams/surgery)
        \item Travel (insurance for travel health, cancellation, etc.)
        \item Dental (covers most cleanings, oral surgery, etc.)
        \item Gender Affirmation Care (certain surgeries and other treatments)
    \end{itemize}

    There is an additional Legal Care Program for \$50 per year (also automatically opted in, billed in student fees). Note that there is different coverage based on whether you take the \textbf{Basic} or \textbf{Enhanced} Care Plan. Cost breakdown can be found \href{https://studentcare.ca/rte/en/McGillUniversitygraduatestudentsPGSS_Cost_HowMuchDoesItCost}{here}.
    \textbf{Keep in mind that students are \textit{automatically} enrolled in ENHANCED Health and Dental Plans!} What can you do in this case?
    \begin{itemize}
        \item Keep the Enhanced plans
        \item Change to Basic plans (you can pick and choose whether you want both Health and Dental to be Basic, or just choose one of them as Basic and the other Enhanced)*
        \item Opt-out of the Health and Dental Plans entirely*
        \item Modify your plan to add family members*
    \end{itemize}

    *For these changes, you MUST do this before the Opt-Out/Changes deadline every semester, which is usually close to the end of September for Fall incoming students, and close to the end of January for Winter incoming students. More details on Opt-outs/changes can be found \href{https://www.studentcare.ca/rte/en/McGillUniversitygraduatestudentsPGSS_ChangeofCoverage_OptOuts}{here}.
    
    Extensive detail surrounding the coverage can be found on PGSS's \href{https://studentcare.ca/rte/en/McGillUniversitygraduatestudentsPGSS_Home}{Studentcare website}. If you want a quick overview, please see this \href{https://pgss.mcgill.ca/document/view/10773/PGSS%20Studentcare%20Health%20and%20Dental%20Insurance%20Presentation.pdf}{document}.

    \subsubsection{Provincial Health Insurance} \index{Health Insurance! Provincial Health Insurance}
    If you are a Canadian from outside of Québec who only came to Montréal to pursue your studies, you are considered a temporary resident in Québec and therefore are not eligible to obtain RAMQ (Régie de l'assurance maladie du Québec). You can be covered by RAMQ if you decide to settle in Québec, please consult \href{https://www.ramq.gouv.qc.ca/en/citizens/health-insurance/register}{the official eligibility check}. \index{Health Insurance!RAMQ}

    If you wish, you can register for RAMQ through this \href{https://www.ramq.gouv.qc.ca/en/citizens/health-insurance/registration-information}{application}. However, it should be noted that anyone aged 26 or older will be required to pay for the Public Prescription Drug Insurance Plan if they do not have access to a private plan. Relevantly, the PGSS Health Insurance is NOT considered a private plan for this purpose. As of Aug 2026, its \href{https://www.ramq.gouv.qc.ca/en/citizens/prescription-drug-insurance/annual-premium}{annual rate} is \$789 per person.
    
    Canadian students from another province can remain covered under their province or territory of origin by virtue of the interprovincial health insurance agreements concluded between the Canadian provinces and territories. To maintain your coverage, you will need to notify your home province of your absences for full-time study purposes. Find the applicable procedure for your home province in \href{https://pgss.mcgill.ca/en/ramq-_-provincial-health-insurance}{this list}. Furthermore, you will need to file taxes in your home province instead of in Québec (see \hyperref[tax]{\textcolor{purble}{\ref*{tax}}}).

    For out-of-province Canadian students, appointments at the McGill's Wellness Hub tend to be the most hassle free as they handle all the insurances behind the scene for you. For off-campus resources, you will have to pay upfront for an appointment, but can submit the receipts to your home province for reimbursements. However, if approved, you are reimbursed according to your province's standard medical fee rates, which may not cover the full amount you were charged in Québec. The difference may then be submitted for reimbursements to the PGSS Health Insurance.

    Canadian non-residents who came here to study have to obtain the International Health Insurance until they have established residency in Canada. These are more specific cases and we encourage students in such situation to reach out directly to \href{https://www.mcgill.ca/student-accounts/tuition-fees/non-tuition-charges/insurance}{McGill Student Services} for help.
    
\subsection{International Students}

    Please refer to Section  
\hyperref[section:inthealth]{\textcolor{purble}{\ref*{section:inthealth}}}
    for more health insurance details specifically applying to international students. The \href{https://www.mcgill.ca/internationalstudents/health/benefits/handbook}{IHI handbook} details exact coverage amounts for each type of health service. 
    
\section{Health Resources}

    \subsection{Physical Health}

    For \textbf{on-campus care}, the \href{https://www.mcgill.ca/wellness-hub/}{Wellness Hub} offers a wide range of services. To meet with a nurse/doctor/counsellor, you can follow \href{https://www.mcgill.ca/wellness-hub/contact}{these instructions} to book an appointment.

    \begin{plaincorollary}
    Appointments with a physician at the Wellness Hub cannot be booked more than a day in advance. Same-day appointments go quickly. You need to call exactly at the opening time to secure a spot. Once on the line, keep pressing 1 instead of waiting for all the prompts.
    \end{plaincorollary}

    For \textbf{off-campus} care, please visit \href{https://www.mcgill.ca/wellness-hub/get-support/find-community-resources/navigatinghealthcare}{this website}, which provides a list of off-campus clinics and healthcare providers. 

    We also get \textbf{Virtual Care}! Our virtual care is through Maple, which provides an easy way to virtually consult physicians on the same day. Depending on the issue, you can get prescribed medication through Maple. More information on Maple and creating an account can be found \href{https://www.mcgill.ca/internationalstudents/health/benefits/maple-virtual-care}{here}.
    

    \subsection{Mental Health}
    For mental health resources, visit the 
    \href{https://www.mcgill.ca/wellness-hub/get-support/mental-health-support}{McGill Wellness Hub}. You can also talk to a licensed-therapist via Maple included in your health insurance. PGSS also provides a list of resources \href{https://pgss.mcgill.ca/en/mental-health-resources}{here} and the physics department has a list of emergency contacts \href{https://www.physics.mcgill.ca/resources/grads.html}{here}. 
    

    \section{How to File Complaints} \index{Complaints}
    

\subsection{Building Complaints}

If you have any issues with the building, you can address your concern to any of the following individuals:
\begin{itemize}
  \item MGAPS Workspace and Accessibility Officer
  \item Physics Building Director: \href{mailto:buildingdirector.physics@mcgill.ca}{buildingdirector.physics@mcgill.ca} (currently Andreas Warburton)
\end{itemize}

\subsection{Safety Complaints}\index{Safety}

For any \textbf{campus safety concerns}, you can:
\begin{itemize}
    \item Call campus security (514-398-3000) or 911. \textbf{These numbers are in case of emergencies}
    \item Call campus security for general inquiries (514-398-4556)
    \item Learn more about campus public safety \href{https://www.mcgill.ca/campussafety/}{here}.
\end{itemize} 

In case of sexual violence/assault, or in case of witnessing sexual violence/assault, you can use these following resources. These resources are meant to be \textbf{entirely confidential}.
\begin{itemize}
    \item Sexual Assault Center of the McGill Students Society (SACOMSS) \href{www.sacomss.org}{www.sacomss.org}
    \item Office for Sexual Violence Response, Support and Education (OVRSE) \href{https://www.mcgill.ca/osvrse/}{www.mcgill.ca/osvrse}
\end{itemize}

\subsection{Supervisor Complaints}

If you have issues with your supervisor that cannot be resolved between you and your supervisor, you next person to contact would be the Graduate Program Director (GPD), who is currently Nik Provatas.
If this is still not working, please visit \href{https://www.mcgill.ca/expmed/policies/conflict-resolution-policy}{this website} for more details on who to contact/how to navigate the hierarchy of help. This is also expanded on in Section \hyperref[section:supercomp]{\textcolor{purble}{\ref*{section:supercomp}}}.

